%======================================================================
\section{투고 서류에 들어가는 것}

\begin{yoji}
PLOS ONE 양식이 요구하는 윤리 진술, 자료 공개, 생성형 AI 사용, 감사의 말, 보조 자료 목록이에요. 영문 원고에서는 윤리부터 생성형 AI까지가 「자료와 방법」 안에, 감사의 말과 보조 자료는 결론 뒤에 있어요.
\end{yoji}

\subsection{윤리 진술}

이 연구는 공개된 블록체인 자료와 공개된 토큰 가격만 썼고, 사람 참여자는 없어요. 점수를 매긴 주소는 계약이에요. 외부 소유 계정은 하나도 점수를 매기거나 보고하지 않았고, 어떤 주소도 실제 신원과 이으려 하지 않았으며, 집계 통계만 보고해요. [윤리 승인 번호는 저자가 넣을 자리예요.]

\subsection{자료 공개}

\begin{itemize}
\item \textbf{저장소의 내용:} 평가 코드, 버전을 붙인 설정 파일, SQL 질의, 분석 계획, W1 등록, 수정 분석을 정의하는 문서, 추출 목록(manifest), 모든 표 뒤에 있는 기계가 읽는 요약과 그 SHA-256 체크섬, 모든 계수 뒤에 있는 계약별 점수와 라벨이에요. 이것들은 Git 저장소 \nolinkurl{https://github.com/abitocodes/phd_works}에서 누구나 받을 수 있어요.
\item \textbf{커밋과 폴더:} 계획(커밋 \texttt{43d8045}), 등록(\texttt{0d3a057}), 등록 분석의 결과(\texttt{8d8b568})는 그 커밋들에서 폴더 \nolinkurl{1-Dissertation-Works/ERC20-Allowance-PageRank-Contract-Reputation/experiments}에 있어요. 수정 분석의 결과는 코드 리비전 \texttt{0a94215}에서 계산했고 커밋 \texttt{950a7a4}로 냈어요. 이 커밋에서는 같은 폴더의 이름이 \nolinkurl{1-Dissertation-Works/ERC20-Allowance-PageRank-Wallet-Reputation/experiments}예요(\url{https://github.com/abitocodes/phd_works/tree/950a7a4cd9401002b3be9348797b4a8dd9d9f075/1-Dissertation-Works/ERC20-Allowance-PageRank-Wallet-Reputation/experiments}).
\item \textbf{요약의 위치:} \nolinkurl{data/3-published-results-for-thesis/}(등록 분석)과 그 아래 폴더 \nolinkurl{revised-usd-12m/}(수정 분석).
\item \textbf{계약별 파일:} 두 분석의 계약별 파일은 공개 Google Drive 폴더에 parquet 파일로 두었어요. 그 링크와 체크섬은 폴더의 파일 \nolinkurl{drive-manifest.json}(커밋 \texttt{ac7d32a})에 있어요.
\item \textbf{대조 스크립트:} 닫힌 식과 트레이더 값을 그 파일들과 대조하는 스크립트는 \nolinkurl{https://github.com/abitocodes/phd-dissertation-draft-en}의 폴더 \nolinkurl{analysis/}에, 이 원고의 표와 그림 1을 만드는 스크립트는 그 저장소의 폴더 \nolinkurl{manuscripts/A-plos-one/}에 있어요.
\item \textbf{보조 자료:} 사전 명세 문서는 S1 File에도 있고, 질의는 S1 Text에 있어요.
\item \textbf{BigQuery:} 공개 BigQuery 표는 버전이 붙은 스냅숏이 아니에요. 그래서 추출한 그래프 표는 저자들의 BigQuery 프로젝트에 두었고, 저장소에는 그 목록이 있어요. 연구 전체의 BigQuery 작업은 6.76~TiB를 청구했고 비용은 USD~42.24였어요.
\item 계약별 파일에는 주소가 공개 장부에 나오는 그대로만 있고, 오프체인 신원은 없어요.
\end{itemize}

\subsection{생성형 AI 사용}

이 연구는 Claude(Anthropic)를 Claude Code 에이전트로 썼고, 언어 다듬기보다 더 많은 일에 썼어요.
\begin{itemize}
\item 첫 번째 저장소에서는 제1저자의 지시에 따라 추출 코드와 평가 코드의 일부를 쓰고 돌렸어요.
\item 두 번째 저장소에서는 폴더 \texttt{analysis/}의 스크립트를 썼어요. 이 스크립트는 공개한 요약에서 닫힌 식과 트레이더 값을 옮겨 오고, 점검을 위해 계약별 파일로 트레이더 값을 다시 계산해요. 이 원고의 표, 그림 1, S1 Table과 S2 Table을 만드는 스크립트도 썼어요.
\item 원고와 보조 자료와 표지 편지의 글을 초안으로 쓰고 고쳤고, 글과 스크립트의 오류를 점검했어요.
\end{itemize}
여기 보고한 모든 수는 스크립트로 된 점검으로 공개한 요약과 대조했어요. 글의 내용에 대한 모든 책임은 저자들에게 있어요.

\subsection{감사의 말}

제1저자의 박사 연구를 지원한 남아프리카 대학교(UNISA) 컴퓨팅 학부에 감사해요. 이 연구는 그 박사 연구의 일부예요. 이벤트 로그는 Google BigQuery에 올라 있는 공개 아비트럼 원 자료에서, 하루 토큰 가격은 공개 DefiLlama API에서 가져왔어요. Claude(Anthropic)를 분석과 글쓰기에 쓴 일은 「자료와 방법」의 생성형 AI 사용에 적었어요.

\subsection{보조 자료}

\begin{itemize}
\item \textbf{S1 Text. 이벤트 풀기, 값 매기기, 라벨 식, 추출 질의.} ERC-20과 GMX~V2 이벤트의 토픽 해시와 풀어낸 필드, 계정 종류 규칙, 토큰 규칙과 달러 값 매기기, 점수 기간과 상수, 스펜더 라벨과 트레이더 라벨의 정확한 식, SQL 틀, 그리고 수로 재지 않은 자료의 성질.
\item \textbf{S1 File. 사전 명세 문서.} 분석 계획(커밋 \texttt{43d8045})과 W1 등록(커밋 \texttt{0d3a057}, SHA-256은 추출 목록에 기록)을 커밋 식별자와 해시와 함께 글자 그대로 옮긴 것.
\item \textbf{S1 Table. 창, 코호트, 그래프, 금액.} 같은 기간 평가와 두 홀드아웃 창의 묘사 통계. 매칭 코호트부터 허락을 받은 트레이더까지.
\item \textbf{S2 Table. 트레이더 결과 비교의 모든 대비.} 두 창, 트레이더 라벨 여섯 개에서 등록한 모든 쌍의 95\% 구간과 97.5\% 구간, 층화한 행. 수정 분석과 등록 분석 모두.
\end{itemize}
